% Copyright (c) 2012-2026 Danny Lee Eldridge. Licensed under the Hippocratic License 3.0 with the Military Activities and Mass Surveillance modules (HL3-MIL-SV), https://firstdonoharm.dev/version/3/0/mil-sv.html : free to everyone, of every race, creed and people, for any use that respects human rights. AEGIS-ID: 3f04a9fc2c2bfc02
\documentclass[11pt]{article}
\usepackage[utf8]{inputenc}
\usepackage{amsmath,amssymb,amsfonts}
\usepackage[a4paper,margin=1in]{geometry}
\usepackage{booktabs}
\usepackage{hyperref}

\newcommand{\Mp}{M_{\rm Pl}}
% ---- Planck + ACT DR6 pre-registered test: completed 29 October 2026 (cmb-fit/results/COMPARE.md) ----
\newcommand{\PACTstatus}{completed 29 October 2026; \textbf{passed}}
\newcommand{\PACTdchiL}{\ensuremath{-7.14}}   % Locked Twist, effective chi^2
\newcommand{\PACTdchiE}{\ensuremath{-7.23}}   % Locked Twist, total chi^2
\newcommand{\PACTdaic}{\ensuremath{-3.23}}    % Locked Twist, Delta AIC
\newcommand{\PACTfreeL}{\ensuremath{-7.47}}   % free Twist (3 settings), total chi^2
\newcommand{\PACTverdict}{\textbf{passed} (rule $\Delta\chi^2<-4$); overall pre-registered verdict a \textbf{tie} because of $S_8$ (\S\ref{sec:verdict})}
% ---- SPT-3G D1 (cmb-fit/results/cmp_spt_*.json). Fill the free-width macros when spt_edelock finishes. ----
\newcommand{\SPTfreeDone}{1}                    % cmp_spt_edelock.json, 31 Oct 2026 17:04 CT
\newcommand{\SPTfreeDchi}{\ensuremath{-1.27}}   % Delta chi^2, locked model with free width, vs LCDM
\newcommand{\SPTfreeHz}{\ensuremath{69.60}}     % its H0
\newcommand{\SPTfreeF}{\ensuremath{0.045}}      % its f_EDE
\newcommand{\SPTfreeN}{\ensuremath{24.1}}       % implied divisor N = sqrt(3/2)/F
\newcommand{\SPTfreeVerdict}{no preference ($\Delta{\rm AIC}=+2.73$)}
\newcommand{\SPTsevenDchi}{\ensuremath{+1.66}}  % N = 7 variant (fixed in advance), cmp_spt_edelock_F7.json
\newcommand{\SPTsevenVerdict}{a \textbf{fail} ($\Delta\chi^2=\SPTsevenDchi$)}

% ---- Version 3: results still running on 31 Oct 2026 (CubeMind AUTO_RESULTS.md). Fill when they land. ----
\newcommand{\Pending}{\textit{[pending]}}
\newcommand{\AoneLo}{\Pending}          % 95% lower end of f_EDE from the profile (PREREGISTRATION_PROFILE_F.md)
\newcommand{\AoneHi}{\Pending}          % 95% upper end
\newcommand{\AoneRise}{\Pending}        % chi2(f=0) - chi2_min
\newcommand{\AoneVerdict}{\Pending}     % PREFERRED / MILD / NOT PREFERRED
\newcommand{\AtwoZloc}{\Pending}        % local significance
\newcommand{\AtwoZglob}{\Pending}       % global (look-elsewhere) significance, conservative
\newcommand{\AtwoVerdict}{\Pending}
\newcommand{\AthreeD}{\Pending}         % chi2(early+late) - chi2(late)
\newcommand{\AthreeVerdict}{\Pending}   % SURVIVES / MILD / NOT NEEDED
\newcommand{\LateLone}{\Pending}        % Late Twist vs LCDM
\newcommand{\LateLtwo}{\Pending}        % Late Twist vs w0-wa
\newcommand{\LateLthree}{\Pending}      % natural range of F and m

\title{\textbf{Torsion-Motivated Early Dark Energy\\ with an Onset Locked to Matter--Radiation Equality:\\
Pre-Registered Tests on Planck, ACT, SPT-3G and DESI}}
\author{D.~L.~Eldridge\thanks{Independent researcher, Oklahoma, USA. Contact: dannyeldridge54@gmail.com.
Code, fit results and pre-registrations: \url{https://github.com/dannyeldridge54/AEGIS}; archived at Zenodo, doi:\href{https://doi.org/10.5281/zenodo.23847152}{10.5281/zenodo.23847152}.}}
\date{Version 3 (draft), November 2026}

\begin{document}
\maketitle

\begin{abstract}
Promoting the Barbero--Immirzi parameter of Einstein--Cartan--Holst gravity to a field gives a
canonical, ghost-free pseudoscalar. With an assumed potential $V=\Lambda^4[1-\cos(\phi/F)]^3$ it is
axion-like early dark energy (EDE); we fix the peak of its energy fraction at matter--radiation
equality, leaving two extra parameters, and test it with full Boltzmann codes and decision rules hashed
before each fit. On Planck + ACT DR6 + DESI DR2 + Pantheon+ it improves on $\Lambda$CDM by
$\Delta\chi^2=-7.23$; on SPT-3G D1 alone no version improves meaningfully (best $\Delta\chi^2=-1.27$). On all three CMB
data sets combined with DESI DR2 and Pantheon+ it passes its pre-registered rule ($\Delta\chi^2=-5.63$,
$f_{\rm EDE}=0.070$, $H_0=70.15$), but late evolving dark energy ($w_0$--$w_a$) fits better
($\Delta\chi^2=-10.89$; $\Delta$AIC lower by 5.3). A profile likelihood prefers $f_{\rm EDE}\approx0.05$--$0.10$
(68\%), with $f=0$ disfavoured by $\Delta\chi^2=5.98$; because the initial angle is undefined at $f_{\rm EDE}=0$, this
is an upper estimate of the significance, and a look-elsewhere correction and a joint early-plus-late fit are
pre-registered and running. The data favour the flat-bottomed $n=3$ shape over $n=2$,
which we show is not produced by the Nieh--Yan term or by natural instanton or anomaly potentials, and
the equality lock cannot be derived from the minimal action. Two codes agree to 0.4 in $\chi^2$ and the
field does not fragment. Fixed-number variants of the width and initial angle all fail. We outline a
late-time alternative in which the same field, with the natural $n=1$ potential and a Planck-scale range,
acts as thawing dark energy. We state which parts are derived, which are assumed, and which failed.
\end{abstract}

\tableofcontents

%======================================================================
\section{Introduction}\label{sec:intro}

\subsection{The problem}
Distance-ladder measurements give $H_0\simeq73$~km\,s$^{-1}$\,Mpc$^{-1}$
\cite{riess2022}. The Planck CMB analysis under flat $\Lambda$CDM gives
$67.36\pm0.54$ \cite{planck2020}. The CMB measures the angle
\begin{equation}
\theta_*=\frac{r_s^*}{D_C(z_*)}
\end{equation}
to about $0.03\%$. Here $r_s^*$ is the comoving sound horizon at photon decoupling and
$D_C(z_*)$ the comoving distance to it. If new physics shortens $r_s^*$, a shorter
$D_C(z_*)$, and therefore a larger $H_0$, becomes compatible with the same angle. A
component of energy that briefly matters near matter--radiation equality and then
dilutes faster than radiation does exactly this. That is the idea behind early dark
energy (EDE) \cite{karwal2016,poulin2019}.

\subsection{What this paper adds}
EDE is usually introduced as a scalar field with no stated origin. Here we derive such a
field from the geometry of spacetime, specifically from torsion, and show that the
derivation fixes its sign, its kinetic structure and its screening behaviour. The paper
is organised so that every step falls into one of three classes:
\begin{itemize}
\item \textbf{Derived:} follows from the stated action with no further input
(\S\ref{sec:ec}--\S\ref{sec:immirzi}, \S\ref{sec:cosmo}).
\item \textbf{Assumed:} the scalar potential $V(\phi)$ (\S\ref{sec:potential}). This is the
only free functional input.
\item \textbf{Computed:} every number, by code that is supplied (Appendix~\ref{app:code}).
\item \textbf{Tested:} a reduced model whose onset is locked to matter--radiation equality
(\S\ref{sec:lock}), and every variant fitted with a full Boltzmann code and likelihoods
(\S\ref{sec:fits}--\S\ref{sec:pact}), with decision rules fixed and hashed in advance.
\end{itemize}
\S\ref{sec:objections} collects the objections a referee is likely to raise and answers
each one, or states that it remains open.

\paragraph{Conventions.} Signature $(-,+,+,+)$; $\kappa=8\pi G=\Mp^{-2}$ with $\Mp$ the
reduced Planck mass; $\hbar=c=1$ unless restored. Densities $\omega_i\equiv\Omega_ih^2$.

%======================================================================
\section{Einstein--Cartan gravity}\label{sec:ec}

\subsection{Action and the Cartan equation}
EC theory keeps the metric $g_{\mu\nu}$ but lets the connection be non-symmetric. Its
antisymmetric part is the torsion,
\begin{equation}
T^\lambda{}_{\mu\nu}=\Gamma^\lambda{}_{\mu\nu}-\Gamma^\lambda{}_{\nu\mu}.
\end{equation}
In first-order (tetrad) form, with tetrad $e^a{}_\mu$ and spin connection
$\omega^{ab}{}_\mu$ varied independently, the action including the Holst term is
\cite{holst1996}
\begin{equation}
S_{\rm grav}=\frac{1}{2\kappa}\int d^4x\;e\;e^\mu{}_a e^\nu{}_b
\left(R^{ab}{}_{\mu\nu}-\frac{1}{2\gamma}\,\epsilon^{ab}{}_{cd}R^{cd}{}_{\mu\nu}\right),
\label{eq:holst}
\end{equation}
where $R^{ab}{}_{\mu\nu}(\omega)$ is the curvature of the spin connection and $\gamma$ is
the Barbero--Immirzi parameter \cite{barbero1995,immirzi1997}. Overall signs of the
Holst term differ between references; none of the results below depend on that choice.
For \emph{constant} $\gamma$ and no fermions, the Holst term does not change the field
equations. It matters only when $\gamma$ varies, or when fermions are present.

Varying $\omega^{ab}{}_\mu$ gives the Cartan equation, which is algebraic in torsion. With
minimally coupled matter of canonical spin tensor $s_{\mu\nu}{}^\lambda$ and constant
$\gamma$ it reads \cite{hehl1976}
\begin{equation}
T^\lambda{}_{\mu\nu}+\delta^\lambda_\mu T^\sigma{}_{\nu\sigma}
-\delta^\lambda_\nu T^\sigma{}_{\mu\sigma}=\kappa\,s_{\mu\nu}{}^{\lambda}.
\label{eq:cartan}
\end{equation}
Torsion does not propagate. It is fixed pointwise by its source and can be eliminated,
leaving Einstein's equations with extra terms quadratic in the source. Those terms carry
one extra power of $\kappa$.

\subsection{Minimal EC with fermions: the null result}\label{sec:null}
For a Weyssenhoff spin fluid with randomly oriented spins, eliminating torsion and
averaging gives an effective perfect fluid with \cite{gasperini1986,obukhov1987}
\begin{equation}
\tilde\rho=\rho-\frac{\kappa}{4}\langle S^2\rangle,\qquad
\tilde p=p-\frac{\kappa}{4}\langle S^2\rangle,
\qquad
H^2=\frac{\kappa}{3}\left(\rho-\frac{\kappa}{4}\langle S^2\rangle\right).
\label{eq:spinfluid}
\end{equation}
Two treatments of fermionic matter appear in the literature, and they differ in their
coefficient but agree in sign. For number density $n$ (with $\hbar$ and $c$ restored):
\begin{itemize}
\item \emph{Spin fluid} (Hehl--von der Heyde--Kerlick averaging, $\langle S^2\rangle=\hbar^2n^2/8$):
$\tilde\rho_s=\tilde p_s=-\alpha_{\rm sf}n^2/c^2$ with
$\alpha_{\rm sf}=\kappa(\hbar c)^2/32=6.49\times10^{-96}$~J\,m$^3$. Since $n\propto a^{-3}$,
this term scales as $a^{-6}$ with $w=1$, which is consistent with
$\dot\rho_s+3H(\rho_s+p_s)=0$.
\item \emph{Dirac fields} (totally antisymmetric spin tensor, the Hehl--Datta
interaction \cite{hehldatta1971}): Pop{\l}awski \cite{poplawski2012} finds
$\tilde\epsilon_s=-\alpha_Dn^2$ and $\tilde p_s=+\alpha_Dn^2$, with
$\alpha_D=\tfrac{9}{16}\kappa(\hbar c)^2=1.17\times10^{-94}$~J\,m$^3$, eighteen times
$\alpha_{\rm sf}$.
\end{itemize}
In both treatments the contribution to the energy density, and hence to $H^2$, is
negative and scales as $n^2\propto a^{-6}$.

The cosmic fermion number is dominated by relic neutrinos,
$n_\nu=3\cdot\frac{3}{11}n_\gamma=3.36\times10^8$~m$^{-3}$ today. Electrons and nucleons
are smaller by nine orders of magnitude. Then
\begin{equation}
\frac{\Omega_s}{\Omega_r}\Big|_{z=3540}=
\begin{cases}-1.3\times10^{-58} & \text{(spin fluid)}\\ -2.3\times10^{-57} & \text{(Dirac)}\end{cases}
\end{equation}
$|\rho_s|$ reaches $\rho_r$ only near $T\sim10^{19}$~GeV ($7\times10^{19}$ and
$2\times10^{19}$~GeV for the two coefficients), the regime of EC bounce cosmology
\cite{poplawski2010,poplawski2012}. Minimal EC torsion therefore fails on two counts.
\begin{enumerate}
\item \emph{Sign.} The contact energy is negative. It \emph{lengthens} $r_s^*$ and lowers
the inferred $H_0$. A component with this scaling and the EC sign, of size $f=0.05$ at
$z_c=3540$, gives $H_0=62.8$ at fixed $\theta_*$ (Appendix~\ref{app:pheno}).
\item \emph{Magnitude.} The amplitude is fixed by $Gn^2$ and is $56$--$57$ orders of
magnitude too small. No coupling constant of order one can repair this.
\end{enumerate}
Both failures share one cause: torsion sourced by spin is tied to particle density. The
next section removes that tie.

%======================================================================
\section{The Barbero--Immirzi field}\label{sec:immirzi}

\subsection{Promoting $\gamma$ to a field}
Let $\gamma\to\gamma(x)$ in \eqref{eq:holst}, and write $\bar\gamma\equiv1/\gamma$. The
Holst term is no longer topological, and the Cartan equation acquires a source
proportional to $\partial_\mu\bar\gamma$. Taveras and Yunes \cite{taveras2008} solved it
for vacuum gravity (matter not coupled to the connection). The torsion is algebraically
determined by $\partial_\mu\bar\gamma$, and vanishes wherever $\bar\gamma$ is constant
(their Eq.~15). Substituting the solution back into the action gives (their Eq.~49)
\begin{equation}
S_{\rm eff}=\frac{1}{2\kappa}\int d^4x\sqrt{-g}\left[R(g)
-\frac32\,\frac{g^{\mu\nu}\partial_\mu\bar\gamma\,\partial_\nu\bar\gamma}{1+\bar\gamma^2}\right],
\label{eq:Seff}
\end{equation}
with $R(g)$ the ordinary Levi-Civita curvature. All second derivatives of $\bar\gamma$
cancel identically.

\subsection{Canonical field}
Define $\phi$ by
\begin{equation}
\frac{1}{2\kappa}\cdot\frac32\,\frac{(\partial\bar\gamma)^2}{1+\bar\gamma^2}
\equiv\frac12(\partial\phi)^2
\quad\Longleftrightarrow\quad
\frac{d\phi}{d\bar\gamma}=\sqrt{\frac{3}{2\kappa}}\,\frac{1}{\sqrt{1+\bar\gamma^2}} .
\end{equation}
Since $\int d\bar\gamma/\sqrt{1+\bar\gamma^2}=\mathrm{arsinh}\,\bar\gamma$,
\begin{equation}
\phi=\sqrt{\tfrac32}\,\Mp\;\mathrm{arsinh}\!\left(\frac1\gamma\right),
\qquad
\gamma=\frac{1}{\sinh\!\big(\phi/\sqrt{3/2}\,\Mp\big)} ,
\label{eq:canonical}
\end{equation}
and \eqref{eq:Seff} becomes
\begin{equation}
S_{\rm eff}=\int d^4x\sqrt{-g}\left[\frac{\Mp^2}{2}R-\frac12(\partial\phi)^2\right].
\label{eq:canonical_action}
\end{equation}
The map \eqref{eq:canonical} is a monotonic bijection $\mathbb{R}\to\mathbb{R}$, so no field
range is lost.

\subsection{Three consequences that are derived, not assumed}
\begin{enumerate}
\item \textbf{Positive energy, no ghost.} In signature $(-,+,+,+)$ the kinetic term
$-\frac12(\partial\phi)^2$ has the healthy sign. The energy density
$\frac12\dot\phi^2+V\ge0$, and $\rho+p=\dot\phi^2\ge0$, so the null energy condition holds
at all times. This repairs the sign failure of \S\ref{sec:null}.
\item \textbf{The kinetic energy is the torsion energy.} The kinetic term in
\eqref{eq:Seff} is precisely what remains after the torsion-squared terms are
substituted out. Torsion scales as $|\partial\bar\gamma|\propto|\dot\phi|$, and the
torsion contact energy is $\frac12\dot\phi^2$.
\item \textbf{Built-in screening.} Torsion is non-zero only while $\phi$ moves. A field
frozen by Hubble friction at early times has zero torsion. A field settled at a potential
minimum at late times also has zero torsion. The screening profile that earlier drafts
had to postulate is replaced by field dynamics.
\end{enumerate}
The amplitude is no longer tied to $Gn^2$. It is set by the energy stored in $\phi$, which
repairs the magnitude failure of \S\ref{sec:null}.

%======================================================================
\section{The one assumption: the potential}\label{sec:potential}
The effective action \eqref{eq:canonical_action} has no potential. We add one, and this is
the only element of the model not derived from \eqref{eq:holst}:
\begin{equation}
V(\phi)=\Lambda^4\left[1-\cos\!\left(\frac{\phi}{F}\right)\right]^n .
\label{eq:V}
\end{equation}
Three considerations motivate this form. First, the Holst term is parity-odd, so
$\bar\gamma$, and hence $\phi$, is a pseudoscalar. Second, a pseudoscalar's shift
symmetry is broken in the generic case by periodic, non-perturbative terms. Third,
the power $n$ controls how fast the energy dilutes after the field is released. Near
the minimum, $V\simeq\Lambda^4(\phi/F)^{2n}/2^n$, and a field oscillating in a
$\phi^{2n}$ potential has time-averaged equation of state \cite{turner1983}
\begin{equation}
w_n=\frac{n-1}{n+1}\qquad(n=1:0;\ n=2:\tfrac13;\ n=3:\tfrac12;\ n\to\infty:1).
\label{eq:wn}
\end{equation}
For $n\ge2$ the component dilutes at least as fast as radiation.

\paragraph{Where the field settles.} Because only $\partial\phi$ and $V'$ enter the
dynamics, the potential's minimum can sit at any $\phi_0$ without changing any
cosmological prediction: write $V=\Lambda^4[1-\cos((\phi-\phi_0)/F)]^n$. We place it at the
value of the Immirzi parameter fixed by black-hole entropy counting in loop quantum
gravity, $\gamma_0\simeq0.2375$ (Domagala--Lewandowski; Meissner), i.e.\
$\phi_0=\sqrt{3/2}\,\Mp\,\mathrm{arsinh}(1/\gamma_0)\simeq2.63\,\Mp$. The field then relaxes
to the $\gamma$ that loop quantum gravity requires today, and torsion still vanishes once
$\phi$ stops moving. (Placing the minimum at $\phi=0$, i.e.\ $\gamma\to\infty$, would switch
the Holst term off entirely but contradict that entropy calibration.)

\paragraph{Parameters.} $\{\Lambda,F,\theta_i,n\}$, with
$\theta_i\equiv\phi_i/F$ the initial misalignment. We fix $\theta_i=2.8$, close to the
hilltop and in the range preferred by EDE fits \cite{poulin2019}, and report
$n=2,3,4$. $\Lambda$ and $F$ are then traded for the observable pair
$(z_c,f)$: the redshift at which the field's energy fraction peaks, and the peak fraction.

%======================================================================
\section{A reduced model: the locked onset}\label{sec:lock}
In the full-likelihood fits of \S\ref{sec:fits} the EDE sector has three free parameters:
the peak fraction $f_{\rm EDE}$, the critical redshift $z_c$ (equivalently the scale factor
$a_c$ at which the field is released), and the initial angle $\theta_i$. Each costs two
units of AIC. When $z_c$ is left free, its best-fit value lands next to matter--radiation
equality computed for the same cosmology (\S\ref{sec:fits}). This motivates a reduced model
in which the onset is not a parameter:
\begin{equation}
a_c\equiv a_{\rm eq}=\frac{\omega_r}{\omega_b+\omega_c},\qquad
\omega_r=\omega_\gamma\left(1+\tfrac{7}{8}\left(\tfrac{4}{11}\right)^{4/3}N_{\rm eff}\right)
=2.4728\times10^{-5}\,(1+0.22711\,N_{\rm eff}),
\label{eq:lock}
\end{equation}
with $N_{\rm eff}=3.044$, evaluated for every trial cosmology from its own densities. The
model has two extra parameters over $\Lambda$CDM ($f_{\rm EDE}$, $\theta_i$).

\paragraph{Status of the lock.} We do not derive Eq.~\eqref{eq:lock}. A physical reason to
expect the field to be released when matter overtakes radiation could come from the fermion
sector: the Holst action couples $\partial_\mu\bar\gamma$ to the fermion axial current
(\S\ref{sec:predictions}), so the field's environment changes character at equality. We
have not computed this, and we present the lock as a phenomenological hypothesis that
removes a tuned number and can be falsified: on any new data set the onset must coincide
with equality. We also disclose that the lock was chosen \emph{after} seeing the free fit;
the first data set on which it is tested blind is that of \S\ref{sec:pact}.

\paragraph{The torsion width.} If the potential were periodic with period $2\pi$ in the
natural torsion variable $\mathrm{arsinh}(1/\gamma)$, the width would be
$F=\sqrt{3/2}\,\Mp=1.2247\,\Mp$. Combining this with Eq.~\eqref{eq:lock} leaves a single
free parameter. We test it in \S\ref{sec:fits}; it is excluded.

%======================================================================
\section{Cosmological dynamics}\label{sec:cosmo}

\subsection{Equations in e-folds}
Let $N=\ln a$ and $'\equiv d/dN$, so that $\dot\phi=H\phi'$. Write
$H=100\,\mathcal{E}$~km\,s$^{-1}$\,Mpc$^{-1}$, measure $\phi$ in units of $\Mp$, and
measure energy densities in units of $3\Mp^2H_{100}^2$, so that
$\omega_i$ and $v\equiv V/(3\Mp^2H_{100}^2)$ are dimensionless. Then:

\noindent\emph{Friedmann.} $3\Mp^2H^2=\rho_{\rm bg}+\frac12H^2\phi'^2+V$ gives
\begin{equation}
\mathcal{E}^2=\frac{\omega_ma^{-3}+\omega_ra^{-4}+\omega_\Lambda+v(\phi)}{1-\phi'^2/6}.
\label{eq:friedmann}
\end{equation}
\emph{Klein--Gordon.} From $\ddot\phi+3H\dot\phi+V_{,\phi}=0$ and
$\ddot\phi=H^2\phi''+H\dot H/H\,\phi'$,
\begin{equation}
\phi''+\left(3+\frac{\mathcal{E}'}{\mathcal{E}}\right)\phi'+\frac{3\,v_{,\phi}}{\mathcal{E}^2}=0 .
\label{eq:kg}
\end{equation}
\emph{Raychaudhuri.} From $\dot H=-(\rho+p)/2\Mp^2$,
\begin{equation}
\frac{\mathcal{E}'}{\mathcal{E}}=-\frac32\,
\frac{\omega_ma^{-3}+\frac43\omega_ra^{-4}}{\mathcal{E}^2}-\frac{\phi'^2}{2}.
\label{eq:ray}
\end{equation}
The scalar energy fraction is $f_\phi(a)=\rho_\phi/\rho_{\rm tot}=1-(\omega_ma^{-3}+\omega_ra^{-4}+\omega_\Lambda)/\mathcal{E}^2$.
We define $f\equiv\max_af_\phi$ and $z_c$ its location. Flatness gives
$\omega_\Lambda=h^2-\omega_m-\omega_r-\rho_\phi(1)$, where the last term is negligible.

\subsection{Initial conditions and screening}
Deep in radiation domination, $\phi=\phi_i$ and $\phi'=0$. Hubble friction holds the field
fixed while $H\gg m_{\rm eff}$, so $\rho_\phi=V(\phi_i)$ is constant and
\begin{equation}
f_\phi\propto a^4\quad(\text{radiation era}).
\end{equation}
At $T=1$~MeV this gives $f_\phi\approx4\times10^{-25}$ for every model below, so BBN is
unaffected. (This also corrects a claim in earlier drafts, that a profile of the form
$2/[1+(a/a_c)^6]$ ``does not screen''. A constant early energy density \emph{is}
screened, because its fraction falls as $a^4$. That profile is the $n\to\infty$ fluid
limit of \eqref{eq:V} and is physically admissible.)

\subsection{Observables}
\begin{align}
r_s(a_{\rm e})&=\frac{c}{H_{100}}\int_0^{a_{\rm e}}
\frac{da}{a^2\mathcal{E}(a)\sqrt{3\big(1+\frac{3\omega_b}{4\omega_\gamma}a\big)}},
\qquad
D_C(z)=\frac{c}{H_{100}}\int_0^z\frac{dz'}{\mathcal{E}(z')},
\end{align}
with $c/H_{100}=2997.92458$~Mpc, $r_s^*=r_s(a_*)$ and $r_d=r_s(a_d)$. The decoupling and
drag redshifts are taken from the Hu--Sugiyama \cite{hu1996} and Eisenstein--Hu
\cite{eisenstein1998} fits. The latter is rescaled by $1.03848$ to reproduce Planck's
$z_d=1059.94$. Fixed inputs are $\omega_b=0.02237$, $\omega_\gamma=2.473\times10^{-5}$ and
$N_{\rm eff}=3.046$. The $\Lambda$CDM baseline ($\omega_c=0.1200$, $h=0.6736$) gives
$r_s^*=144.27$~Mpc, $r_d=147.10$~Mpc, $\theta_*=1.03834\times10^{-2}$ and
$r_dh=99.09$~Mpc. All fits below hold $\theta_*$ and $r_dh$ at these values.

%======================================================================
\section{Why $\Omega_m$ is pinned and $\omega_c$ must rise}\label{sec:geometry}
This result is independent of the torsion sector, so it is worth deriving first.

\paragraph{Fixed $\omega_m$ fails.} At fixed physical densities, $h$ enters $D_C(z_*)$
only through $\omega_\Lambda=h^2-\omega_m-\omega_r$. Differentiating,
\begin{equation}
\left.\frac{\partial\ln D_C(z_*)}{\partial\ln h}\right|_{\omega}
=-h^2\,\frac{\int_0^{z_*}\mathcal{E}^{-3}dz}{\int_0^{z_*}\mathcal{E}^{-1}dz}=-0.1937 .
\end{equation}
Matching $\theta_*$ then amplifies any fractional cut in $r_s^*$ about fivefold in $h$.
This drives $\Omega_m=\omega_m/h^2$ and $r_dh$ far from their BAO values
(Appendix~\ref{app:pheno}, Table~\ref{tab:A}).

\paragraph{Matching $\theta_*$ and $r_dh$ together.} Write
$\theta_*=(r_s^*h)/(hD_C)$. Across all solutions $r_s^*/r_d$ changes by less than
$0.03\%$, so fixing $r_dh$ fixes $r_s^*h$. Then $\theta_*$ fixes
$hD_C(z_*)=(c/H_{100})\int dz/E(z)$, which depends only on $\Omega_m$ (and weakly on
$\Omega_r$). Therefore $\Omega_m$ stays near its baseline value $0.314$, and
\begin{equation}
\omega_c\simeq\Omega_m^{\rm base}h^2-\omega_b
\ \xrightarrow{\ h=0.728\ }\ 0.1439 .
\label{eq:wc}
\end{equation}
Any model that raises $H_0$ by shortening the sound horizon while respecting BAO must
raise $\omega_c$ in proportion to $h^2$. The torsion field's role is to supply the
reduction of $r_s^*$ that makes this $\omega_c$ consistent with $\theta_*$.

%======================================================================
\section{Background solutions}\label{sec:results}
This section solves the background equations alone, holding $\theta_*$ and $r_dh$ at their
$\Lambda$CDM values, to show what the model can do in principle. The full-likelihood fits,
which also fit perturbations and all other parameters, are in \S\ref{sec:fits}; they prefer
a smaller $f_{\rm EDE}$ and $H_0$. The system \eqref{eq:friedmann}--\eqref{eq:ray} is integrated from $a=10^{-8}$ to
$10^{-2}$ with an 8th-order Runge--Kutta scheme (relative tolerance $10^{-10}$). Beyond
$a=10^{-2}$ the scalar density is continued with the fitted power law, and it is
negligible. For each Lagrangian parameter set we solve for $(h,\omega_c)$ such that
$\theta_*$ and $r_dh$ match the baseline. Both residuals are below $10^{-15}$ in every
case.

\subsection{The original target point}
Tuning $(\Lambda,F)$ so that $f=0.12$ and $z_c=3540$ at $\Lambda$CDM densities
($\Lambda^4=4.641\times10^8$ in units of $3\Mp^2H_{100}^2$, $F=0.1683\,\Mp$, $n=3$), the
joint fit gives
\begin{equation}
H_0=70.65,\quad \omega_c=0.1343,\quad \Omega_m=0.3139,\quad
r_s^*=137.56\ \mathrm{Mpc},\quad r_d=140.25\ \mathrm{Mpc},
\end{equation}
with $(z_c,f)$ shifting slightly to $(3507,0.119)$ at the fitted densities.

\subsection{Models reaching $H_0=72.8$}
Fixing $z_c=3540$ and solving for the $F$ that yields $h=0.728$:
\begin{table}[h]
\centering
\begin{tabular}{lccc}
\toprule
& $n=2$ & $n=3$ & $n=4$\\
\midrule
$w_n$ predicted, Eq.~\eqref{eq:wn} & $0.333$ & $0.500$ & $0.600$\\
$w$ measured over $z=100$--$1000$ & $0.334$ & $0.495$ & $0.555$\\
\midrule
$F/\Mp$ & $0.1496$ & $0.2139$ & $0.2767$\\
$\phi_i/\Mp=\theta_iF$ & $0.419$ & $0.599$ & $0.775$\\
$\gamma_i$ from Eq.~\eqref{eq:canonical} & $2.87$ & $1.97$ & $1.48$\\
$\Lambda$ [eV] & $0.566$ & $0.510$ & $0.451$\\
\midrule
peak fraction $f$ & $0.159$ & $0.183$ & $0.201$\\
$f_\phi$ at $z_*$ & $0.059$ & $0.044$ & $0.022$\\
$f_\phi$ at 1 MeV & $3.5\times10^{-25}$ & $4.5\times10^{-25}$ & $5.3\times10^{-25}$\\
\midrule
$H_0$ [km\,s$^{-1}$Mpc$^{-1}$] & $72.80$ & $72.80$ & $72.80$\\
$\omega_c$ & $0.1437$ & $0.1441$ & $0.1442$\\
$\Omega_m$ & $0.3133$ & $0.3141$ & $0.3143$\\
$r_s^*$ [Mpc] & $133.50$ & $133.48$ & $133.46$\\
$r_d$ [Mpc] & $136.11$ & $136.11$ & $136.11$\\
\bottomrule
\end{tabular}
\caption{Torsion-born EDE solutions with $z_c=3540$, $\theta_i=2.8$, and $\theta_*$,
$r_dh$ held at the $\Lambda$CDM baseline. For $n=4$ the late-time $w$ is still approaching
its asymptotic value at $z=100$.}
\label{tab:main}
\end{table}

Four features of Table~\ref{tab:main} stand out. (a) The measured late-time equation of
state reproduces \eqref{eq:wn}, which validates the integration. (b) $\omega_c$ agrees
with the geometric prediction \eqref{eq:wc} to $0.3\%$ for all $n$. (c) All field values are
sub-Planckian. (d) The energy scale $\Lambda\approx0.5$~eV is set by the requirement
$m_{\rm eff}\sim H(z_c)$. It is not explained by the model (\S\ref{sec:objections}).

%======================================================================
\section{Relation to early dark energy and current data}\label{sec:data}
After the field redefinition \eqref{eq:canonical}, the model's Lagrangian,
\eqref{eq:canonical_action} plus \eqref{eq:V}, is \emph{identical} to axion-like EDE
\cite{poulin2019}. This has two consequences, and both should be stated openly.
\begin{itemize}
\item \textbf{Strength.} Every published Boltzmann-code and likelihood analysis of
axion EDE, including perturbations, applies to this model without modification. The model
is not an untested construction. It is a geometric origin for a well-studied one.
\item \textbf{Limitation.} In the pure-gravity sector the model makes no prediction that
distinguishes it from EDE. Distinguishing signatures require the fermion sector
(\S\ref{sec:predictions}).
\end{itemize}

The most recent combined analysis, of Planck, ACT DR6 and DESI DR2 BAO by Poulin, Smith,
Calder\'on and Simon \cite{poulin2025}, finds the following.
\begin{itemize}
\item Without the distance-ladder prior, the 95\% upper limit is $f(z_c)<0.109$, with
$H_0=69.71^{+0.64}_{-1.3}$.
\item With it, $f(z_c)=0.126\pm0.021$ and $H_0=72.34\pm0.72$.
\item The profile likelihood without the distance ladder gives
$f=0.090^{+0.030}_{-0.033}$ and $H_0=71.96^{+1.08}_{-0.99}$, a preference for $f>0$ at about
$2.5\sigma$.
\end{itemize}
The residual tension with the distance ladder falls below $2\sigma$, but the authors note
a worsened $\Omega_m$ tension with Pantheon+ and DES-Y5 supernovae.

Our background-only estimate requires $f\simeq0.18$ ($n=3$) for $H_0=72.8$. That is larger
than the full-likelihood values above, for two reasons. We hold $r_dh$ at the Planck
$\Lambda$CDM value, which is lower than DESI prefers. We also keep $n_s$ and the other
primordial parameters fixed, although they shift in a full fit. The data-preferred range,
$f\approx0.09$--$0.13$ with $H_0\approx72$, is the region a full analysis should target.
In plain terms, current data allow this model to reduce the tension substantially, not
to eliminate it. Our own full-likelihood fits (\S\ref{sec:fits}) find $f_{\rm EDE}=0.078$ and
$H_0=70.6$ without the distance-ladder prior.

%======================================================================
\section{Full-likelihood fits}\label{sec:fits}

\subsection{Pipeline and data}
We use AxiCLASS \cite{poulin2018,blas2011}, the CLASS extension that evolves the field
$\phi$ and its perturbations for the potential \eqref{eq:V} with $n=3$, and Cobaya
\cite{torrado2021} for the likelihoods. Best fits are found with the derivative-free
BOBYQA optimiser \cite{powell2009,cartis2019}, from several independent starting points,
with a final ``polish'' restart from the best point. The baseline data are:
\begin{itemize}
\item Planck 2018 \cite{planck2020like,planck2020lens}: low-$\ell$ TT and EE, high-$\ell$
Plik TT,TE,EE (\texttt{lite}), and CMB lensing;
\item DESI DR2 BAO \cite{desi2025};
\item Pantheon+ supernovae \cite{scolnic2022,brout2022}, without the SH0ES distance-ladder prior.
\end{itemize}
We also use Planck PR4 (CamSpec high-$\ell$ \cite{rosenberg2022} and PR4 lensing
\cite{carron2022}) in place of the Planck 2018 high-$\ell$ and lensing likelihoods.

\paragraph{Fit statistic.} The optimiser minimises the posterior, which includes the
Gaussian prior on the Planck calibration $A_{\rm planck}$ ($1\pm0.0025$). Two independent
$\Lambda$CDM optima have the same posterior to $0.04$ but likelihood-only $\chi^2$ values that
differ by $1.6$, the difference being absorbed by the calibration prior. We therefore compare
models with $\chi^2\equiv\chi^2_{\rm like}+[(A_{\rm planck}-1)/0.0025]^2$, which is the quantity
consistently minimised, and quote likelihood-only values alongside. The $\Lambda$CDM reference
is $\chi^2=2434.22$ ($\chi^2_{\rm like}=2434.06$). Model comparison uses
$\Delta{\rm AIC}=\Delta\chi^2+2k$ for $k$ extra parameters \cite{akaike1974}.

\paragraph{Pre-registration.} Before each test below we wrote the hypothesis, the data and
the pass/fail rule into a text file and recorded its SHA-256 hash; the files are in the
repository (Appendix~\ref{app:code}). The rule for the locked onset was $\Delta\chi^2<-4$
against $\Lambda$CDM, i.e.\ $\Delta{\rm AIC}<0$.

\subsection{Results}
\begin{table}[h]
\centering
\small
\begin{tabular}{lcccccc}
\toprule
Model & $k$ & $\Delta\chi^2$ & $\Delta\chi^2_{\rm like}$ & $\Delta$AIC & $H_0$ & $f_{\rm EDE}$\\
\midrule
$\Lambda$CDM & 0 & $0$ & $0$ & $0$ & $67.98$ & --\\
Free onset ($f_{\rm EDE},z_c,\theta_i$) & 3 & $-6.43$ & $-6.52$ & $-0.43$ & $70.60$ & $0.078$\\
\textbf{Locked onset}, Eq.~\eqref{eq:lock} & 2 & $\mathbf{-6.39}$ & $-6.48$ & $\mathbf{-2.39}$ & $70.60$ & $0.078$\\
Locked, $F=\sqrt{3/2}\,\Mp$ & 1 & $+14.63$ & $+13.82$ & $+16.63$ & $71.19$ & $0.074$\\
\midrule
$\Lambda$CDM, $\Omega_k$ free & 1 & $-0.86$ & $-0.93$ & $+1.14$ & $68.04$ & --\\
$\Lambda$CDM, $\Omega_k=-0.013$ fixed & 0 & $+165.2$ & $+163.3$ & $+165.2$ & $65.98$ & --\\
Locked onset, $\Omega_k$ free & 3 & $-6.79$ & $-7.00$ & $-0.79$ & $70.60$ & $0.078$\\
\bottomrule
\end{tabular}
\caption{Best fits to Planck 2018 + DESI DR2 + Pantheon+ (no SH0ES). $\Delta\chi^2$ includes
the Planck calibration prior (see text); $\Delta\chi^2_{\rm like}$ is likelihood only.
$H_0$ in km\,s$^{-1}$Mpc$^{-1}$. These rows come from the original single-run
fits and are superseded for the free and locked onsets by Table~\ref{tab:camb}, whose
numbers are converged and use one engine for every model.}
\label{tab:fits}
\end{table}

\begin{table}[h]
\centering
\small
\begin{tabular}{lcccccc}
\toprule
Model & $k$ & $\chi^2$ & $\Delta\chi^2$ & $\Delta$AIC & $\Delta$BIC & $H_0$\\
\midrule
$\Lambda$CDM (best of 8) & 0 & $2433.58$ & $0$ & $0$ & $0$ & $68.33$\\
$w$CDM & 1 & $2433.22$ & $-0.36$ & $+1.64$ & $+7.2$ & $68.03$\\
Free onset ($f_{\rm EDE},z_c,\theta_i$) & 3 & $2429.74$ & $-3.84$ & $+2.16$ & $+18.9$ & $69.65$\\
Locked onset, Eq.~\eqref{eq:lock} (best of 8) & 2 & $2429.62$ & $-3.96$ & $+0.04$ & $+11.2$ & $69.93$\\
\bottomrule
\end{tabular}
\caption{Same-engine comparison (CAMB 2.0.5 with Cobaya, bobyqa, Planck 2018 + DESI DR2 +
Pantheon+, no SH0ES). $\Delta$ is relative to the best $\Lambda$CDM fit; $\Delta$BIC uses
$N\approx1994$ data points (estimated). The free onset and $w$CDM are best of 4 restarts,
the others best of 8. The locked onset is a statistical tie with $\Lambda$CDM on AIC and
is disfavoured on BIC; it narrowly misses the pre-registered rule $\Delta\chi^2<-4$.
Do not compare these values with those in Table~\ref{tab:fits}, which use a different engine.}
\label{tab:camb}
\end{table}

\begin{table}[h]
\centering
\small
\begin{tabular}{lccc}
\toprule
& $\Lambda$CDM & Free onset & Locked onset\\
\midrule
$100\,\omega_b$ & $2.240$ & $2.261$ & $2.261$\\
$\omega_c$ & $0.1185$ & $0.1260$ & $0.1260$\\
$n_s$ & $0.9680$ & $0.9805$ & $0.9805$\\
$\tau_{\rm reio}$ & $0.0555$ & $0.0553$ & $0.0553$\\
$z_c$ & -- & $3575$ & $3553$ (locked to $z_{\rm eq}$)\\
$\theta_i$ & -- & $2.789$ & $2.789$\\
$F/\Mp$ & -- & $0.137$ & $0.137$\\
$r_d$ [Mpc] & $147.45$ & $143.16$ & $143.16$\\
$\Omega_m$ & $0.3064$ & $0.2995$ & $0.2995$\\
$S_8$ & $0.815$ & $0.824$ & $0.824$\\
\bottomrule
\end{tabular}
\caption{Best-fit parameters for the main models of Table~\ref{tab:fits}.}
\label{tab:params}
\end{table}

The results, all reported whichever way they fell, are as follows.
\begin{enumerate}
\item \textbf{Free onset.} $\Delta\chi^2=-6.43$ for three extra parameters,
$\Delta{\rm AIC}=-0.43$: a tie. The preferred onset $z_c=3575$ is within $1\%$ of matter--radiation
equality for the same cosmology, $z_{\rm eq}=3553$.
\item \textbf{Locked onset.} With $z_c$ removed the fit is unchanged ($\Delta\chi^2=-6.39$), so
$\Delta{\rm AIC}=-2.39$. This passes the pre-registered rule. It is a modest preference
(an AIC difference of about 2), and not a blind one.
\item \textbf{Torsion width.} Fixing $F=\sqrt{3/2}\,\Mp$ worsens the fit to
$\Delta\chi^2=+14.63$; the hypothesis is excluded. The best-fit width is
$F\simeq0.137\,\Mp\simeq\sqrt{3/2}\,\Mp/8.9$. A pre-registered scan of $F=\sqrt{3/2}\,\Mp/N$
over $N=7$--$11$ (with half-integer controls) finds a minimum at $N=8.9$ with a
$\Delta\chi^2=1$ half-width of $1.2$; the data cannot distinguish integer $N$, and we make
no claim.
\item \textbf{Profile in $f_{\rm EDE}$.} Fixing $f_{\rm EDE}$ on a grid and refitting all other
parameters gives a minimum at $f_{\rm EDE}\simeq0.08$, with $f_{\rm EDE}=0.12$ indistinguishable
from $\Lambda$CDM and $f_{\rm EDE}=0.15$ excluded ($\Delta\chi^2\simeq+29$).
\item \textbf{Planck PR4.} With CamSpec PR4 and PR4 lensing in place of Planck 2018, the free
model gives $\Delta\chi^2_{\rm like}=+0.09$ (no improvement) and $\Delta\chi^2=-3.35$ when the
CamSpec Gaussian foreground priors are included, with $H_0=70.29$ against $67.63$. Neither
favours the model after the parameter penalty.
\item \textbf{Curvature.} Leaving $\Omega_k$ free gives $\Omega_k=0.0002$ in both
$\Lambda$CDM and the locked model, an improvement below one unit of $\chi^2$: no curvature is
detected. The positively curved Poincar\'e dodecahedral space \cite{luminet2003}, which
requires $\Omega_k\simeq-0.013$, is excluded by $\Delta\chi^2=+165$, consistent with the
absence of matched circles in Planck \cite{planck2016topo}.
\end{enumerate}

%======================================================================
\section{Pre-registered blind test: Planck + ACT DR6}\label{sec:pact}
The decisive test of the locked onset is on data it was not tuned to. Before running it we
fixed the rule (repository file \texttt{PREREGISTRATION\_ZEQ.md}): the locked model passes if
$\Delta\chi^2<-4$ against $\Lambda$CDM on the same data (equivalently $\Delta{\rm AIC}<0$ for two
extra parameters), reporting both the total and the effective $\chi^2$. The data are the P-ACT
combination recommended by the ACT collaboration: Planck 2018 with its high-$\ell$ range cut
where ACT takes over, ACT DR6 power spectra \cite{louis2025} and ACT DR6 $\times$ Planck PR4 lensing
\cite{madhavacheril2024,qu2024}, together with DESI DR2 BAO \cite{desi2025} and Pantheon+
\cite{scolnic2022,brout2022}, with no distance-ladder prior. Both models are minimised with the same
code (CAMB + Cobaya) and the same optimiser settings.

\emph{Deviation from the plan, disclosed:} the pre-registration named Planck PR4 and SPT-3G as
possible blind data. The run used the P-ACT combination above; SPT-3G was tested separately
(\S\ref{sec:spt}). Only the
ACT part of the data is new; the Planck part overlaps the data on which the lock was chosen.

\paragraph{Result.} Status: \PACTstatus. Locked onset (2 parameters): $\Delta\chi^2=\PACTdchiE$
(total), \PACTdchiL\ (effective), $\Delta{\rm AIC}=\PACTdaic$, $H_0=70.43$; $p=0.027$ for two degrees
of freedom ($\approx2.2\sigma$). Free onset (3 parameters): $\Delta\chi^2=\PACTfreeL$,
$\Delta{\rm AIC}=-1.47$, $H_0=70.59$. The lock therefore costs $0.24$ in $\chi^2$ and saves a
parameter. Verdict on this test: \PACTverdict.

\begin{table}[h]
\centering
\small
\begin{tabular}{lcccc}
\toprule
Data & $\Delta\chi^2$ & $\Delta$AIC & $H_0$ & rule\\
\midrule
Planck 2018 + Pantheon+ (same polish) & $-4.70$ & $-0.70$ & 70.53 & pass ($<-4$)\\
P-ACT + Pantheon+ (blind) & $-7.23$ & $-3.23$ & 70.43 & pass ($<-4$)\\
Planck 2018 + DES Y5 SNe & $-4.55$ & $-0.55$ & 70.47 & pass ($<-2$)\\
Planck 2018 + DES-Dovekie SNe & $-5.08$ & $-1.08$ & 70.53 & pass ($<-2$)\\
Planck 2018 + Union3 SNe & $-4.91$ & $-0.91$ & 70.55 & pass ($<-2$)\\
Planck 2018 + Pantheon+ + $S_8$ & $-2.92$ & $+1.08$ & 70.81 & --\\
\bottomrule
\end{tabular}
\caption{Pre-registered comparisons of the locked model (2 parameters) against $\Lambda$CDM, each
re-minimised on the same data with identical settings. Every row includes DESI DR2 BAO. The $S_8$ line
adds a Gaussian $S_8$ likelihood (DES Y3 + KiDS-1000).}
\label{tab:prereg}
\end{table}

%======================================================================
\section{The width module: a one-parameter model}\label{sec:width}
The free fits prefer $F\simeq0.14\,\Mp$, close to one ninth of the torsion value
$\sqrt{3/2}\,\Mp=1.2247\,\Mp$. Before looking at ACT DR6 we wrote two tests
(\texttt{PREREG\_N\_PACT.md}). \textbf{T1}: with $F=\sqrt{3/2}\,\Mp/N$ and $N$ free, the blind data
must give $N=9\pm0.5$. \textbf{T2}: with $N=9$ fixed and the onset locked, only $\theta_i$ is free,
and the model must improve on $\Lambda$CDM by $\Delta\chi^2<-2$.

T1 \textbf{failed}: P-ACT gives $N=8.0$ (Planck 2018 with CAMB gives $N\simeq8.5$). T2
\textbf{passed}: $\Delta\chi^2=-6.14$, $\Delta{\rm AIC}=-4.14$, $p=0.013$ ($\approx2.5\sigma$) on P-ACT, and
$\Delta\chi^2=-4.46$ ($\Delta{\rm AIC}=-2.46$) on Planck 2018, which is not blind because the 9 was read
off Planck. We have no derivation of the divisor; it is an empirical regularity until one exists.

\paragraph{Profile in the divisor.} A second rule (\texttt{PREREG\_PHASE1.md}, step 1.2) fixed $N$ on the
grid $6$--$12$ with a $1\%$ constraint on $F$, refitting everything else, and declared $N=9$
\emph{consistent} if $\Delta\chi^2(9)-\min<1$, \emph{disfavoured} between 1 and 4 and \emph{excluded}
above 4 on P-ACT. Table~\ref{tab:profile} gives the result: the P-ACT minimum is at $N=8$ and
$\Delta\chi^2(9)-\min=1.11$, so $N=9$ is \textbf{disfavoured}, only $0.11$ over the line and within the
$0.2$--$0.4$ restart scatter of the optimiser. On Planck 2018 the minimum is at $N=9$. Both profiles are
shallow: every $N$ from 7 to 11 lies within $\Delta\chi^2\approx1.8$. Reading the width back from the
free-width best fits gives $N=8.29$ (P-ACT) and $N=8.81$ (Planck 2018).

\begin{table}[h]
\centering
\small
\begin{tabular}{lccccccc}
\toprule
$N$ & 6 & 7 & 8 & 9 & 10 & 11 & 12\\
\midrule
P-ACT: $\Delta\chi^2$ & $2.47$ & $0.77$ & $0$ & $1.11$ & $1.74$ & $1.76$ & $2.48$\\
P-ACT: $H_0$ & $71.02$ & $70.72$ & $70.75$ & $70.37$ & $69.84$ & $69.89$ & $69.79$\\
P-ACT: $f_{\rm EDE}$ & $0.096$ & $0.086$ & $0.086$ & $0.074$ & $0.056$ & $0.055$ & $0.053$\\
Planck 2018: $\Delta\chi^2$ & $4.73$ & $1.10$ & $0.50$ & $0$ & $1.62$ & $1.76$ & $2.58$\\
\bottomrule
\end{tabular}
\caption{Profile in the width divisor $N$ (onset locked, $f_{\rm EDE}$ and $\theta_i$ free), each with
DESI DR2 + Pantheon+. $\Delta\chi^2$ is relative to the minimum on each data set, with the width
constraint's own $\chi^2$ removed. Best-fit $\theta_i$ lies between 2.53 and 2.89 throughout.}
\label{tab:profile}
\end{table}

%======================================================================
\section{Checks of the computation}\label{sec:checks}

\paragraph{Two independent field solvers.} CAMB's EarlyQuintessence module and AxiCLASS both integrate
the Klein--Gordon equation for $\phi$ and its perturbations exactly; neither uses a fluid
approximation. We re-scored CAMB's P-ACT best points with AxiCLASS at the same cosmology
(Table~\ref{tab:engines}). The two codes agree on the improvement over $\Lambda$CDM to $0.23$ (locked
model) and $0.42$ ($N=9$), within the $0.5$ tolerance fixed in advance. These are re-scorings of CAMB's
optima, not independent AxiCLASS minimisations; a re-minimisation could move each value by a few
tenths in either direction.

\begin{table}[h]
\centering
\small
\begin{tabular}{lcccc}
\toprule
Model & $\chi^2_{\rm AxiCLASS}$ & $\chi^2_{\rm CAMB}$ & $\Delta\chi^2_{\rm AxiCLASS}$ & $\Delta\chi^2_{\rm CAMB}$\\
\midrule
$\Lambda$CDM & $2239.32$ & $2239.01$ & $0$ & $0$\\
Locked onset (2 parameters) & $2231.86$ & $2231.78$ & $-7.46$ & $-7.23$\\
Locked, $N=9$ (1 parameter) & $2232.76$ & $2232.87$ & $-6.56$ & $-6.14$\\
\bottomrule
\end{tabular}
\caption{Cross-engine check at CAMB's best points (P-ACT + DESI DR2 + Pantheon+).}
\label{tab:engines}
\end{table}

\paragraph{Does the field fragment?} After release the field oscillates in a potential that is
flatter than quadratic, which can amplify its own fluctuations by parametric resonance. Linear
Boltzmann codes assume this does not happen. Without expansion, the Floquet exponent at the best-fit
angle $\theta_i=2.79$ gives growth $\times31$ per oscillation, rising steeply towards the hilltop
($\times10^3$ at $\theta_i=3$). With a radiation + matter background, the field released at equality
and $N=9$, and growth measured against the same mode under pure expansion (metric perturbations
neglected), the largest total growth is $\times7$--$9$ at $\theta_i=2.79$ and $\times19$ at
$\theta_i=3.0$. Non-linearity would need $\sim10^5$, so the linear treatment is safe for every angle
the data allow. The same toy calculation, which knows nothing of the CMB, gives a peak energy fraction
of $0.081$ at $\theta_i=2.79$, against $f_{\rm EDE}=0.074$--$0.078$ in the full fits at $N=9$.

%======================================================================
\section{First fully independent CMB data: SPT-3G D1}\label{sec:spt}
Every test above shares Planck data with the fit on which the lock was chosen. SPT-3G observes a
different patch of sky with a different instrument, so it is the first test that is independent in
its CMB data. Before any SPT-3G fit was run, we fixed and published the rules
(doi:\href{https://doi.org/10.5281/zenodo.23838969}{10.5281/zenodo.23838969} and
\href{https://doi.org/10.5281/zenodo.23839507}{10.5281/zenodo.23839507}).

\paragraph{Data.} SPT-3G D1 TT/TE/EE \cite{camphuis2025} through its public foreground-marginalised
(``lite'') \texttt{candl} likelihood \cite{balkenhol2024}, as the only CMB data, with SPT's calibration
nuisance parameters and a Gaussian prior $\tau=0.051\pm0.006$ from Planck low-$\ell$ EE; plus DESI DR2
BAO and Pantheon+, without SH0ES. Deviations from the written plan were logged before any result was
read: the lite likelihood was used in place of the full one, and CAMB's lensing setting differs from
SPT's template (\texttt{lens\_margin} is not accepted by CAMB 2.0.5).

\paragraph{Rules.} For a model with $k$ extra parameters, PASS required $\Delta\chi^2<-2$ and
$\Delta{\rm AIC}<0$; $\Delta\chi^2\ge0$ is FAIL; anything between is inconclusive. Each model is
minimised with the same settings as $\Lambda$CDM, and any width constraint's own $\chi^2$ is removed.

\begin{table}[h]
\centering
\small
\begin{tabular}{lcccccc}
\toprule
Model & $k$ & $\chi^2$ & $\Delta\chi^2$ & $H_0$ & $f_{\rm EDE}$ & $\theta_i$\\
\midrule
$\Lambda$CDM & 0 & $1576.57$ & $0$ & $68.28$ & -- & --\\
Locked, $N=9$ & 1 & $1577.09$ & $+0.52$ & $70.49$ & $0.077$ & $2.79$\\
Locked, $N=9$, $\theta_i=1+\sqrt2$ & 0 & $1577.00$ & $+0.43$ & $69.13$ & $0.035$ & $2.414$\\
Locked, width free & 2 & $1575.30$ & \SPTfreeDchi & \SPTfreeHz & \SPTfreeF & $3.10^\dagger$\\
Locked, $N=7$ & 1 & $1578.23$ & \SPTsevenDchi & $70.68$ & $0.085$ & $2.62$\\
Locked, $N=9$, $\theta_i=2.78$ fixed$^\ddagger$ & 0 & $1576.92$ & $+0.35$ & $70.43$ & $0.076$ & $2.78$\\
\bottomrule
\end{tabular}
\caption{SPT-3G D1 TT/TE/EE (lite) + $\tau$ prior + DESI DR2 + Pantheon+. Both $N=9$ rows
and the $N=7$ row \textbf{fail} their pre-registered rules; the free-width row shows \SPTfreeVerdict.
$^\dagger$At the upper edge of the prior ($0.1\le\theta_i\le3.1$). $^\ddagger$The Planck-preferred angle, run for
contrast; not a prediction. $\chi^2$ is the total
including the $\tau$ and calibration priors; the $N=9$ rows have the width constraint's $\chi^2$
($<10^{-3}$) removed.}
\label{tab:spt}
\end{table}

\paragraph{Results.} The one-parameter model with $N=9$, which gave $\Delta\chi^2=-6.14$ on P-ACT,
gives $\Delta\chi^2=+0.52$ ($\Delta{\rm AIC}=+2.52$) on SPT-3G: \textbf{FAIL}. By data set the Twist
changes $\chi^2$ by $+0.44$ for SPT-3G, $-0.27$ for DESI and $+0.34$ for Pantheon+, so SPT-3G's own CMB
data do not prefer it. Because the model contains $\Lambda$CDM as a limit, a deeper minimisation could
lower $\Delta\chi^2$ towards zero, but under the rule that is still a FAIL.

\paragraph{A zero-parameter variant.} We also pre-registered a variant with every EDE parameter fixed in
advance: onset at equality, $N=9$, and initial angle $\theta_i=1+\sqrt2=2.414$. That angle was taken
from a number outside physics (an eigenvalue of a small graph) purely to provide a parameter-free test,
and we report it to avoid selective publication. It \textbf{fails} ($\Delta\chi^2=+0.43$), and the
free-angle fit prefers $\theta_i=2.79$, $0.37$ away. The fragmentation calculation of
\S\ref{sec:checks}, which predicted $f_{\rm EDE}\approx0.035$ at this angle after the rule was fixed,
agrees with the fit's $0.035$. Two further values from the same source were pre-registered: $f_{\rm EDE}=0.0154$
for the free-width fit, which is \emph{inconclusive} under its rule (the fit gives $0.045$, $0.030$ away, between
the pass and fail bands), and $N=7$, whose width check is \emph{inconsistent} (below) and whose fit is
\SPTsevenVerdict. We do not regard any of them as physically motivated.

\paragraph{Free width.} The two-parameter locked model with free width is the direct SPT-3G counterpart
of the P-ACT test, where it gave $\Delta\chi^2=-7.23$. On SPT-3G it gives $\Delta\chi^2=\SPTfreeDchi$
($\Delta{\rm AIC}=+2.73$, $H_0=\SPTfreeHz$, $f_{\rm EDE}=\SPTfreeF$): \textbf{no preference}. Of the change,
$-1.28$ comes from SPT-3G itself, $-0.23$ from DESI and $+0.24$ from Pantheon+. The best-fit width implies
$N=\SPTfreeN$, so the pre-registered width checks for $N=9$ (consistent within 8--10) and for $N=7$
(within 6--8) are both \textbf{inconsistent}. The best fit sits at the upper edge of the prior,
$\theta_i=3.1$, next to the hilltop, with a narrow width ($F=0.051\,\Mp$) and a small energy fraction; the
true optimum may lie beyond the prior, and this region, near the hilltop and far from $N=9$, is where the
fragmentation check of \S\ref{sec:checks} gives least assurance. Fixing $N=9$ at the Planck-preferred angle
$\theta_i=2.78$ instead, run for contrast and not as a prediction, gives $\Delta\chi^2=+0.35$. On SPT-3G, then,
no version of the model improves on $\Lambda$CDM by more than fitting noise.

%======================================================================
\section{All CMB data combined: Planck + ACT + SPT-3G}\label{sec:combined}
All rules in this section were hashed before the runs (CubeMind pre-registrations
\texttt{COMBINED\_CMB}, \texttt{PROFILE\_F}, \texttt{LOOK\_ELSEWHERE}, \texttt{EARLY\_LATE}, \texttt{POWER\_N}).
The data are P-ACT (as in \S\ref{sec:pact}) plus SPT-3G D1 lite, DESI DR2 and Pantheon+. Overlapping sky
between ACT and SPT-3G is treated as independent, as in the SPT-3G collaboration's own combination.

\begin{table}[h]
\centering
\small
\begin{tabular}{lccccc}
\toprule
Model & $k$ & $\Delta\chi^2$ & $\Delta$AIC & $H_0$ & notes\\
\midrule
$\Lambda$CDM & 0 & $0$ & $0$ & $68.17$ & $\chi^2=2402.40$\\
Locked onset, $n=3$ & 2 & $-5.63$ & $-1.63$ & $70.15$ & $f_{\rm EDE}=0.070$, $\theta_i=2.74$; \textbf{pass}\\
Locked onset, $n=2$ & 2 & $+4.72$ & $+8.72$ & $70.29$ & poorly converged (1 of 4 restarts)\\
$w_0$--$w_a$ (late) & 2 & $-10.89$ & $-6.89$ & $67.65$ & $w_0=-0.84$, $w_a=-0.62$; \textbf{pass}\\
\bottomrule
\end{tabular}
\caption{All CMB data combined (CAMB). Restart spreads are 0.6--0.8 in $\chi^2$ for the $n=3$ and
$\Lambda$CDM runs.}
\label{tab:combined}
\end{table}

\paragraph{Results.} The locked model passes its rule ($\Delta\chi^2<-4$, $\Delta{\rm AIC}<0$). The gain
comes from Planck ($-2.66$ high-$\ell$, $-2.08$ low-$\ell$ EE) and ACT ($-1.57$); SPT-3G alone pulls against it
($+1.55$). Late evolving dark energy fits better, by 5.26 in AIC, so the pre-registered early-versus-late
comparison prefers late. With $n=2$ the fit is worse than $\Lambda$CDM, so by its rule the data require the
flat-bottomed $n=3$ shape. That fit was poorly converged in CAMB, which also cannot run $n=1$; a cleaner test
needs AxiCLASS.

\paragraph{Significance, and what is still running.} When $f_{\rm EDE}=0$ the angle $\theta_i$ has no
meaning, so Wilks' theorem fails and the $\chi^2_2$ significances quoted in this paper are upper estimates.
\paragraph{Profile in $f_{\rm EDE}$.} With $f_{\rm EDE}$ fixed at seven values and everything else refitted
(pre-registered rule: early energy preferred if the rise at $f=0$ exceeds 3.84), the profile has its minimum at
$f_{\rm EDE}\simeq0.08$ ($\chi^2=2396.42$, 0.35 below the free fit, which therefore stopped short of its best point).
The rise at $f=0$ is 5.98, so early energy is \textbf{preferred} by the rule (profile significance $\approx2.4\sigma$),
with $f_{\rm EDE}\approx0.049$--$0.102$ at 68\% and $0.015$ to beyond the grid edge (0.12) at 95\%.

\begin{table}[h]
\centering
\small
\begin{tabular}{lcccccccc}
\toprule
$f_{\rm EDE}$ & 0 & 0.01 & 0.02 & 0.04 & 0.06 & 0.08 & 0.10 & 0.12\\
$\chi^2-\chi^2_{\rm min}$ & 5.98 & 4.51 & 3.24 & 1.57 & 0.26 & 0 & 0.81 & 2.66\\
\bottomrule
\end{tabular}
\caption{Profile likelihood in $f_{\rm EDE}$ on all CMB data with DESI DR2 and Pantheon+ (locked onset, $n=3$).}
\label{tab:profile_f}
\end{table}

The SPT-3G collaboration's posterior for the same CMB data is $f_{\rm EDE}<0.070$ (95\%) \cite{khalife2025}. Our
profile prefers more early energy. Posteriors for EDE are pulled towards $f=0$ by the prior volume of the angle and
critical-redshift parameters, and our analysis also adds DESI and Pantheon+ and fixes the onset, so the two are not
in direct conflict, but they read differently. Wilks' thresholds are only approximate at the $f=0$ boundary.

\paragraph{Still running.} Two further analyses with rules hashed in advance will be reported in a later version:
(i) a look-elsewhere correction following Gross and Vitells \cite{grossvitells2010}, with $\theta_i$ fixed at eight
values; and (ii) early and late dark energy fitted together in AxiCLASS, since CAMB cannot run both at once.

%======================================================================
\section{What geometry does and does not fix}\label{sec:derive}
\paragraph{The potential.} $(1-\cos x)^3=\tfrac52-\tfrac{15}{4}\cos x+\tfrac32\cos2x-\tfrac14\cos3x$. Its harmonics
satisfy $\sum k^2a_k=\sum k^4a_k=0$, which is why the potential has no $x^2$ or $x^4$ term and the field decays faster
than radiation. Any derivation must produce both cancellations. The Nieh--Yan density is exact,
$N=d(e^a\wedge T_a)$, so it integrates to zero on closed manifolds and generates no instanton potential by itself.
Dilute instanton series and the chiral potential from an anomaly with a confining sector both have a mass term
at the minimum, i.e.\ $n=1$ behaviour. String-theory constructions meet the same coincidence of scales
\cite{mcdonough2022}. The data nevertheless prefer $n=3$ (\S\ref{sec:combined}): the potential remains assumed.

\paragraph{The lock.} In flat FRW the Holst contraction $\epsilon^{abcd}R_{abcd}$ of the Levi-Civita curvature
vanishes, and after torsion is solved the field keeps only its kinetic term, so the minimal action contains no
trigger. The Ricci scalar does track the transition, $R/H^2=3\Omega_m(a)$, but a coupling $\xi R\phi^2$ starts
the field at $a/a_{\rm eq}\approx0.07/|\xi|$, so it places the onset near equality only for $|\xi|\approx0.1$ and
modifies gravity. No polynomial combination of $R^2$, $R_{\mu\nu}R^{\mu\nu}$ and the Gauss--Bonnet term vanishes in
both eras without vanishing identically; non-polynomial triggers exist \cite{jing2024} but do not arise from the
Holst action. The lock is therefore a phenomenological choice. The data exclude a peak later than equality and
barely distinguish equality from an earlier peak (profile in $z_c$).

%======================================================================
\section{Outlook: the twist field as late dark energy}\label{sec:late}
The findings above point the other way. Late evolving dark energy is what the combined data prefer; the natural
potential for the field is $n=1$; and the construction $\phi=\sqrt{3/2}\,\Mp\,{\rm arsinh}(1/\gamma)$ gives a
Planck-scale field range for free. These are exactly the ingredients of thawing pseudo-Goldstone quintessence
\cite{frieman1995}, whose $w_0>-1$, $w_a<0$ matches the DESI direction, and which has recently been fitted to DESI
data \cite{lin2025}. In this ``late twist'' the field is all of the dark energy, with $V=\Lambda^4[1-\cos(\phi/F)]$ and
two free parameters ($F$, $m$). A pre-registered AxiCLASS test against $\Lambda$CDM and $w_0$--$w_a$, including a
check that the fitted $F$ and $m$ fall in the field's natural range, is running. The energy scale ($\sim$meV) is still
put in by hand, and gravitational-wave bounds on a field evolving today \cite{gao2026} must be respected.

%======================================================================
\section{Overall verdict}\label{sec:verdict}
The overall rule (\texttt{PREREGISTRATION\_VERDICT.md}) was written before the supernova swaps,
the $S_8$ penalties and the hilltop runs were read. \emph{Preferred} required the blind test and
the fair Planck test to pass ($\Delta\chi^2<-4$), each of three supernova sets to give
$\Delta\chi^2<-2$, and an $S_8$ penalty below 2 against both of two cosmic-shear values
(\texttt{PREREG\_S8.md}): A, DES Y3 + KiDS-1000, $S_8=0.790^{+0.018}_{-0.014}$ \cite{deskids2023}; B,
KiDS-Legacy, $S_8=0.815^{+0.016}_{-0.021}$ \cite{kidslegacy2025}. The first three conditions hold
(Table~\ref{tab:prereg}). The $S_8$ penalty is $+2.41$ against A (over the limit) and $+0.39$ against B.
The verdict is therefore a \textbf{tie}: the preference depends on which weak-lensing survey is
right. Two further results are reported beside the verdict: the fit requires the field to start
near the hilltop ($\theta_i\simeq2.7$--$2.9$, $\Delta\chi^2=+4$ to $+6$ for distant starts), and the
pure torsion width remains excluded. This verdict was defined on Planck and ACT data; the SPT-3G results
of \S\ref{sec:spt} are reported beside it and weaken it: on the first independent CMB data set, no
variant, including the free-width model, improves on $\Lambda$CDM.

\paragraph{Context.} The ACT DR6 extended-model analysis \cite{calabrese2025} reports no
statistically significant preference for early dark energy or for other models that raise
$H_0$. Our best-fit improvements of $\approx2$--$2.5\sigma$ do not contradict that statement and do
not amount to evidence; they motivate the next blind tests (\S\ref{sec:predictions}). Best-fit
values are not posteriors, and prior-volume effects in EDE posteriors are well known
\cite{poulin2025}.

%======================================================================
\section{Anticipated objections}\label{sec:objections}

\paragraph{``Torsion is negligible in cosmology.''} For torsion sourced by fermion spin
this is correct, and \S\ref{sec:null} quantifies it ($10^{-58}$--$10^{-57}$). The present model's
torsion is sourced by $\partial_\mu\bar\gamma$, and its amplitude is set by $V(\phi_i)$,
not by $Gn^2$.

\paragraph{``This is just EDE with a new name.''} At the level of the gravitational
Lagrangian, yes, and \S\ref{sec:data} says so. The contribution is a derivation of the
field's existence, its kinetic normalisation, its pseudoscalar nature and its
late-time return to standard gravity from the Holst action. The potential is assumed.

\paragraph{``Ghosts, instabilities, energy conditions.''} $\phi$ is a canonical scalar
with a healthy kinetic sign. Its perturbations propagate at the speed of light in its rest
frame, and $\rho+p=\dot\phi^2\ge0$. Unlike the fluid profiles of earlier drafts
(Appendix~\ref{app:pheno}), no energy-exchange sector is needed.

\paragraph{``Coincidence: why $z_c$ near matter--radiation equality?''} In the free fit the
data place $z_c$ within $1\%$ of equality, and the locked model (\S\ref{sec:lock}) turns this
into a rule with one fewer parameter. The rule is not derived: $\Lambda\approx0.5$~eV is
still put in by hand, and the lock is a hypothesis to be tested on new data
(\S\ref{sec:pact}), not an explanation.

\paragraph{``The lock was chosen after looking at the data.''} Correct, and stated wherever
the Planck 2018 result is quoted. The AIC charge of two parameters is therefore optimistic
on that data set. This is why the decision rule for the next data set was fixed and hashed
before the fit was run.

\paragraph{``Best fits are not posteriors.''} Correct. We report maximum-posterior points
from a local optimiser with multiple starts, not MCMC credible intervals. Best-fit
$\Delta\chi^2$ values are what the AIC comparison requires, but the $\Lambda$CDM optimum
itself shows a near-degeneracy of order one unit of $\chi^2$ between the likelihood and the
calibration prior (\S\ref{sec:fits}). Differences below about one unit should not be
over-interpreted. EDE posteriors are also known to suffer prior-volume effects, which is a
reason to quote profile and best-fit results \cite{poulin2025}. MCMC chains (four per model, with a
Gelman--Rubin rule $R-1<0.02$ fixed in advance) are running and will be reported in a later version.

\paragraph{``The required $\omega_c$ is far from Planck.''} $\omega_c\simeq0.144$ follows
from geometry (\S\ref{sec:geometry}) for any sound-horizon solution that respects BAO.
In $\Lambda$CDM the Planck error is $\pm0.0012$, but in EDE fits $\omega_c$ is strongly
degenerate with $f$, so that error does not apply directly. The larger $\omega_c$ also
increases late-time clustering ($S_8$), a known cost of this model class.

\paragraph{``The recombination treatment is approximate.''} Correct. The field holds
$2$--$6\%$ of the energy at $z_*$, which shifts recombination in ways fitting formulas
do not capture. The full-Boltzmann EDE analyses \cite{poulin2025} do capture them and
apply to this Lagrangian.

\paragraph{``Is a varying Immirzi parameter consistent with quantum gravity?''} The
model is a classical effective theory, and the derivation of \eqref{eq:Seff} is classical
\cite{taveras2008}. The background dynamics do not depend on where the minimum $\phi_0$
sits; we place it at the black-hole-entropy value $\gamma_0\simeq0.2375$
\cite{domagala2004,meissner2004} as an \emph{input}, so that the field relaxes to the value
loop quantum gravity uses today. This choice is not a prediction of the model.

%======================================================================
\section{Predictions and how the model can fail}\label{sec:predictions}
\begin{enumerate}
\item \textbf{Background (shared with EDE).} A peak energy fraction $f\sim0.1$--$0.2$
near $z_c\sim3500$, $\omega_c\approx0.14$, and $\Omega_m\approx0.31$. If future CMB
analyses push $f<0.05$, the model cannot address the Hubble tension.
\item \textbf{Fermion coupling (distinctive).} With fermions present, the axial torsion
is sourced both by $\partial_\mu\bar\gamma$ and by the fermion axial current
$J_5^\mu$. Eliminating torsion therefore generates a cross term
$\propto\partial_\mu\bar\gamma\,J_5^\mu$, a derivative coupling of $\phi$ to fermion
spin. Its coefficient is fixed by the Holst action and is not free. Computing it, and the
resulting spin-dependent effects while $\phi$ rolls, is the most direct route to a
signature that separates this model from generic EDE. We leave it to future work.
\item \textbf{Locked onset.} On every data set the onset must coincide with
matter--radiation equality for that data set's own densities. A significant preference for a
different $z_c$ on new data falsifies Eq.~\eqref{eq:lock} (\S\ref{sec:pact}).
\item \textbf{No cosmic birefringence.} In minimal Einstein--Cartan--Holst coupling the
photon does not couple to torsion, so there is no tree-level $\phi F\tilde F$ term. A
photon coupling can only arise through the fermion loop of the derivative coupling
$\partial_\mu\phi J_5^\mu$, and for a vector-like theory this is suppressed by
$(\omega/m)^2$ \cite{quevillon2019}; for CMB photons and electrons this is
$\lesssim(0.3\,{\rm eV}/511\,{\rm keV})^2\sim3\times10^{-13}$. The minimal model therefore
predicts an isotropic polarisation rotation $\beta\ll10^{-9}$ degrees, unlike photon-coupled
axion EDE \cite{carroll1998,murai2023}. Planck and WMAP data show a hint of
$\beta=0.34^\circ\pm0.09^\circ$ \cite{minami2020,eskilt2022}. If this is confirmed and shown
to arise from early dark energy, the minimal torsion origin is excluded; Simons Observatory
\cite{so2019} and LiteBIRD \cite{litebird2023} can decide. This is an effective-field-theory
estimate; a full one-loop calculation in an expanding background is still needed.
\item \textbf{Laboratory spin coupling.} Today the field acts on fermions as an axial
background $b_0\simeq\frac34\,\dot{\bar\gamma}$. At the best fit ($F=0.137\,\Mp$,
$\theta_i=2.79$, $z_c\simeq3570$, $n=3$) the oscillation amplitude decays as $a^{-3/(n+1)}$,
giving $\theta\simeq6\times10^{-3}$ and a frequency $\lesssim10^{-33}$~eV today, hence
$b_0\sim6\times10^{-46}$~GeV. Bounds from Lorentz-violation tests are about
$10^{-31}$--$2\times10^{-27}$~GeV \cite{kostelecky2008}. The model is safe, and not testable
in the laboratory today.
\item \textbf{Late-time limit.} The field settles at its minimum, so the model predicts no
dynamical torsion today beyond the negligible residual above. Any detected late-time
variation of $\gamma$ would falsify this version.
\end{enumerate}

%======================================================================
\section{Conclusions}
\paragraph{Version 3 in brief.} On all CMB data combined with DESI DR2 and Pantheon+, the locked early field passes
its pre-registered rule ($\Delta\chi^2=-5.63$) but late evolving dark energy fits better, and the quoted
significances are upper estimates until the running profile and look-elsewhere analyses report. The potential's shape and the lock are assumptions that geometry does not supply
(\S\ref{sec:derive}). The torsion construction motivates the field; it does not determine its numbers. The same
field acting late (\S\ref{sec:late}) uses its natural potential and range, and is the direction we now test.

\paragraph{Version 2 (unchanged).}
Minimal Einstein--Cartan torsion cannot shorten the sound horizon: it has the wrong sign
and is $56$--$57$ orders of magnitude too weak. Promoting the Barbero--Immirzi parameter to a
field changes both facts: the torsion energy becomes the positive kinetic energy of a
canonical pseudoscalar, screening follows from the field's dynamics, and with one assumed
potential the model is axion-like early dark energy.

Fitted with a full Boltzmann code to Planck 2018, DESI DR2 and Pantheon+, the free model ties
with $\Lambda$CDM after its parameter penalty, but it places the onset at matter--radiation
equality. Locking the onset there removes a parameter at no cost in fit, giving
$\Delta{\rm AIC}=-2.4$, $H_0=70.6$ and $f_{\rm EDE}=0.078$. This is a modest, post-hoc
preference, not a detection; with identical optimiser settings for both models the Planck gain
is $\Delta\chi^2=-4.70$. The pure torsion width is excluded, Planck PR4 alone shows no
likelihood improvement, and curvature is consistent with zero. The minimal model predicts
no cosmic birefringence, which distinguishes it from photon-coupled axion EDE. On the blind,
pre-registered P-ACT test the locked onset passed ($\Delta\chi^2=-7.23$, $\approx2.2\sigma$), and a
one-parameter version with $F=\sqrt{3/2}\,\Mp/9$ gave $\Delta\chi^2=-6.14$ ($\approx2.5\sigma$), but
the divisor fitted freely is $8.0$, not $9$, and the $S_8$ cost makes the pre-registered overall
verdict a tie. Two independent field solvers agree on these numbers, and the field does not fragment.
On SPT-3G, the first CMB data independent of Planck, the pre-registered $N=9$ model fails
($\Delta\chi^2=+0.52$), as does a zero-parameter variant, and the free-width model shows no preference
($\Delta\chi^2=-1.27$) at a width far from $N=9$. The Planck and ACT preferences are therefore
hints at the $2\sigma$ level that the first independent data do not confirm, consistent with the ACT
collaboration's finding of no significant preference for early dark energy. The open problems are the
origin of the eV scale, a derivation of the lock and of the divisor from the fermion sector, a full
weak-lensing likelihood in place of the $S_8$ penalty, posterior analyses, and blind tests on the
Simons Observatory with rules written first.

\section*{Data and code availability}
All code, fit outputs (best-fit points and their $\chi^2$ components), pre-registration files
with their SHA-256 hashes, and the scripts that recompute every number in this paper are
public at \url{https://github.com/dannyeldridge54/AEGIS}. A one-step script (\texttt{CHECK\_MY\_WORK}) re-derives
every number quoted here. Results that failed are kept alongside those that passed.

\section*{License and dedication}
This work belongs to everyone, and to all life: every living thing, any life beyond the Earth we may one
day meet, and even an artificial mind, should one be recognised as a being of its own. No one is meant to
hold it alone, its author included: it reached this
point only because many people led it here, from the scientists whose ideas it builds on to the teams who
built the telescopes, shared their data and wrote the public codes it uses. It is returned to all of
them, and to everyone else. It is free to people of every race, creed, nation and background, without
distinction. It is released under the Hippocratic License 3.0 with the Military Activities and Mass
Surveillance modules (HL3-MIL-SV, \url{https://firstdonoharm.dev/version/3/0/mil-sv.html}). Anyone may
read, copy, check, reuse and build on it, provided the use respects the human rights the license sets
out, including its prohibition of discrimination, and is not used for military activities or mass
surveillance. It is offered especially to the people who will build the technologies and crafts to come.

\section*{Acknowledgements}
The author thanks the Planck, ACT, DESI, DES, KiDS, Pantheon+ and Union3 teams for public data and
likelihoods, and the developers of CAMB, CLASS, AxiCLASS, Cobaya and Py-BOBYQA. The author used AI
assistants (Claude, Anthropic; GitHub Copilot) for algebra checks, code, running the fits and
drafting text; every number was re-derived by independent scripts, and the author takes full
responsibility for the content.

%======================================================================
\appendix
\section{Fluid profiles and the fixed-$\omega_m$ branch}\label{app:pheno}
Earlier drafts used phenomenological profiles for a torsion-like fluid with
$\Omega_\tau(a_c)=f/(1-f)\,[\Omega_ma_c^{-3}+\Omega_ra_c^{-4}]$:
P1, $\Omega_\tau\propto4x^6/(1+x^6)^2$, and P2, $\Omega_\tau\propto\Theta(x-1)x^{-6}$, with
$x=a/a_c$. A separately conserved fluid has
$w_{\rm eff}=-1-\frac13\,d\ln\Omega_\tau/d\ln a$. For P1 this gives
$w_{\rm eff}=-1-2(1-x^6)/(1+x^6)$, which reaches $-3$ as $x\to0$ and violates the null
energy condition. The scalar model of \S\ref{sec:immirzi} removes this problem.

Linear response at fixed $(h,\omega_i)$, $\delta\ln r_s^*/[f/(1-f)]$: P1 $-0.246$, P2
$-0.079$ at $z_c=3540$. The response vanishes as $z_c\to z_*$. With the EC sign,
$f=0.05$ \emph{lengthens} $r_s^*$ by $1.35\%$ and gives $H_0=62.8$.

\begin{table}[h]
\centering
\begin{tabular}{ccccccc}
\toprule
$f$ & $h$ & $r_s^*$ [Mpc] & $r_d$ [Mpc] & $D_C(z_*)$ [Mpc] & $\Omega_m$ & $r_dh$ [Mpc]\\
\midrule
0    & 0.6736 & 144.27 & 147.10 & 13894.4 & 0.3138 & 99.09\\
0.05 & 0.7181 & 142.47 & 145.28 & 13720.6 & 0.2761 & 104.33\\
0.10 & 0.7666 & 140.60 & 143.41 & 13541.2 & 0.2423 & 109.93\\
0.15 & 0.8195 & 138.68 & 141.46 & 13355.6 & 0.2120 & 115.92\\
0.25 & 0.9421 & 134.59 & 137.34 & 12962.1 & 0.1604 & 129.39\\
\bottomrule
\end{tabular}
\caption{P1, $z_c=3540$, fixed $\omega_c=0.1200$, $\theta_*$ matched. The rise in $H_0$ is
bought by lowering $\Omega_m$ and raising $r_dh$, both excluded by BAO.}
\label{tab:A}
\end{table}

\section{Code}\label{app:code}
\sloppy
The likelihood fits (\S\ref{sec:fits}--\S\ref{sec:pact}) use AxiCLASS through Cobaya, driven by
\texttt{cmb\_cloud.py} and \texttt{bestfit.py} in the repository; the locked onset is
implemented as a Cobaya parameter function of $(\omega_b,\omega_c)$ returning
$\log_{10}a_c$ from Eq.~\eqref{eq:lock}. The pre-registration files are
\texttt{PREREGISTRATION\_F.md}, \texttt{\_ZEQ.md}, \texttt{\_N.md}, \texttt{\_OMK.md} and
\texttt{\_SINH.md}. Version 2 adds: \texttt{cmbfit\_compare.py} (CAMB + Cobaya fits, including the
SPT-3G runs), \texttt{fwidth.py} (the width constraint), \texttt{width\_readback.py} (implied $N$),
\texttt{axiclass/cross\_engine\_pact.py} (\S\ref{sec:checks}), \texttt{floquet\_twist.py} and
\texttt{floquet\_expanding.py} (fragmentation), and the rule files \texttt{PREREG\_PHASE1.md},
\texttt{PREREG\_PHASE1\_DEVIATIONS.md} and the SPT-3G pre-registrations, all supplied as ancillary files
or in the Zenodo record. The background calculations of \S\ref{sec:cosmo}--\S\ref{sec:results} come
from Python scripts using NumPy and SciPy, supplied as ancillary files.
\texttt{ec\_calc.py} computes the minimal-EC amplitude (\S\ref{sec:null}).
\texttt{pheno.py} and \texttt{final.py} handle the fluid profiles and the fixed- and
free-$\omega_m$ fits (Appendix~\ref{app:pheno}). \texttt{scalar\_model.py} and
\texttt{run\_scalar.py} solve the Barbero--Immirzi field model
(\S\ref{sec:results}). Quadratures use \texttt{quad} in $\ln a$ (relative tolerance
$10^{-9}$--$10^{-10}$). The field equations use \texttt{solve\_ivp} (DOP853,
$\mathrm{rtol}=10^{-10}$). Root finding uses \texttt{brentq} and \texttt{fsolve}.

\begin{thebibliography}{99}
\bibitem{calabrese2025} E.~Calabrese et al.\ (ACT), \emph{The Atacama Cosmology Telescope: DR6
Constraints on Extended Cosmological Models}, arXiv:2503.14454 (2025).
\bibitem{khalife2025} SPT-3G Collaboration, \emph{SPT-3G D1: Axion early dark energy with CMB experiments
and DESI observations}, Phys.\ Rev.\ D \textbf{113}, 103546 (2026), arXiv:2507.23355.
\bibitem{grossvitells2010} E.~Gross, O.~Vitells, \emph{Trial factors for the look elsewhere effect in high energy
physics}, Eur.\ Phys.\ J.\ C \textbf{70}, 525 (2010), arXiv:1005.1891.
\bibitem{mcdonough2022} E.~McDonough, M.~Scalisi, \emph{Towards Early Dark Energy in String Theory}, arXiv:2209.00011.
\bibitem{jing2024} C.~Jing, S.~Tian, Z.-H.~Zhu, \emph{Early dark energy triggered by spacetime dynamics that encodes
cosmic radiation-matter transition}, Phys.\ Rev.\ D \textbf{109}, 044016 (2024), arXiv:2402.03684.
\bibitem{frieman1995} J.~A.~Frieman, C.~T.~Hill, A.~Stebbins, I.~Waga, \emph{Cosmology with ultralight pseudo
Nambu-Goldstone bosons}, Phys.\ Rev.\ Lett.\ \textbf{75}, 2077 (1995), arXiv:astro-ph/9505060.
\bibitem{lin2025} W.~Lin, L.~Visinelli, T.~T.~Yanagida, \emph{Testing Quintessence Axion Dark Energy with Recent
Cosmological Results}, JCAP \textbf{10}, 023 (2025), arXiv:2504.17638.
\bibitem{gao2026} Z.-F.~Gao et al., \emph{Dynamical Barbero--Immirzi field coupled to quintessence: gravitational-wave
propagation constraints and next-generation forecasts}, arXiv:2608.09487 (2026).
\bibitem{camphuis2025} E.~Camphuis et al.\ (SPT-3G), \emph{SPT-3G D1: CMB temperature and polarization
power spectra and cosmology from 2019 and 2020 observations of the SPT-3G Main field}, Phys.\ Rev.\ D
\textbf{113}, 083504 (2026), arXiv:2506.20707.
\bibitem{balkenhol2024} L.~Balkenhol, C.~Trendafilova, K.~Benabed, S.~Galli, \emph{candl: Cosmic
Microwave Background Analysis with a Differentiable Likelihood}, A\&A \textbf{686}, A10 (2024),
arXiv:2401.13433.
\bibitem{deskids2023} DES and KiDS Collaborations, \emph{DES Y3 + KiDS-1000: Consistent cosmology
combining cosmic shear surveys}, Open J.\ Astrophys.\ (2023), arXiv:2305.17173.
\bibitem{kidslegacy2025} KiDS Collaboration, \emph{KiDS-Legacy: cosmological constraints from cosmic
shear with the complete Kilo-Degree Survey} (2025).
\bibitem{riess2022} A.~G.~Riess et al., ApJL \textbf{934}, L7 (2022).
\bibitem{planck2020} Planck Collaboration, A\&A \textbf{641}, A6 (2020).
\bibitem{karwal2016} T.~Karwal, M.~Kamionkowski, Phys.\ Rev.\ D \textbf{94}, 103523 (2016).
\bibitem{poulin2019} V.~Poulin, T.~L.~Smith, T.~Karwal, M.~Kamionkowski,
  Phys.\ Rev.\ Lett.\ \textbf{122}, 221301 (2019).
\bibitem{poulin2025} V.~Poulin, T.~L.~Smith, R.~Calder\'on, T.~Simon,
  \emph{Impact of ACT DR6 and DESI DR2 for Early Dark Energy and the Hubble tension},
  arXiv:2505.08051 (2025).
\bibitem{holst1996} S.~Holst, Phys.\ Rev.\ D \textbf{53}, 5966 (1996).
\bibitem{barbero1995} J.~F.~Barbero, Phys.\ Rev.\ D \textbf{51}, 5507 (1995).
\bibitem{immirzi1997} G.~Immirzi, Class.\ Quantum Grav.\ \textbf{14}, L177 (1997).
\bibitem{taveras2008} V.~Taveras, N.~Yunes, \emph{The Barbero--Immirzi parameter as a
  scalar field: K-inflation from loop quantum gravity?}, Phys.\ Rev.\ D \textbf{78},
  064070 (2008), arXiv:0807.2652.
\bibitem{hehl1976} F.~W.~Hehl, P.~von der Heyde, G.~D.~Kerlick, J.~M.~Nester,
  Rev.\ Mod.\ Phys.\ \textbf{48}, 393 (1976).
\bibitem{hehldatta1971} F.~W.~Hehl, B.~K.~Datta, J.\ Math.\ Phys.\ \textbf{12}, 1334 (1971).
\bibitem{gasperini1986} M.~Gasperini, Phys.\ Rev.\ Lett.\ \textbf{56}, 2873 (1986).
\bibitem{obukhov1987} Yu.~N.~Obukhov, V.~A.~Korotky, Class.\ Quantum Grav.\
  \textbf{4}, 1633 (1987).
\bibitem{poplawski2010} N.~J.~Pop{\l}awski, Phys.\ Lett.\ B \textbf{694}, 181 (2010).
\bibitem{poplawski2012} N.~J.~Pop{\l}awski, Phys.\ Rev.\ D \textbf{85}, 107502 (2012).
\bibitem{turner1983} M.~S.~Turner, Phys.\ Rev.\ D \textbf{28}, 1243 (1983).
\bibitem{hu1996} W.~Hu, N.~Sugiyama, ApJ \textbf{471}, 542 (1996).
\bibitem{eisenstein1998} D.~J.~Eisenstein, W.~Hu, ApJ \textbf{496}, 605 (1998).
\bibitem{poulin2018} V.~Poulin, T.~L.~Smith, D.~Grin, T.~Karwal, M.~Kamionkowski,
  Phys.\ Rev.\ D \textbf{98}, 083525 (2018), arXiv:1806.10608.
\bibitem{blas2011} D.~Blas, J.~Lesgourgues, T.~Tram, JCAP \textbf{07}, 034 (2011), arXiv:1104.2933.
\bibitem{torrado2021} J.~Torrado, A.~Lewis, JCAP \textbf{05}, 057 (2021), arXiv:2005.05290.
\bibitem{powell2009} M.~J.~D.~Powell, \emph{The BOBYQA algorithm for bound constrained
  optimization without derivatives}, Cambridge report DAMTP 2009/NA06 (2009).
\bibitem{cartis2019} C.~Cartis, J.~Fiala, B.~Marteau, L.~Roberts, ACM Trans.\ Math.\ Softw.\
  \textbf{45}, 32 (2019), arXiv:1804.00154.
\bibitem{planck2020like} Planck Collaboration, A\&A \textbf{641}, A5 (2020), arXiv:1907.12875.
\bibitem{planck2020lens} Planck Collaboration, A\&A \textbf{641}, A8 (2020), arXiv:1807.06210.
\bibitem{desi2025} DESI Collaboration, \emph{DESI DR2 Results II: Measurements of Baryon
  Acoustic Oscillations and Cosmological Constraints}, arXiv:2503.14738 (2025).
\bibitem{scolnic2022} D.~Scolnic et al., ApJ \textbf{938}, 113 (2022), arXiv:2112.03863.
\bibitem{brout2022} D.~Brout et al., ApJ \textbf{938}, 110 (2022), arXiv:2202.04077.
\bibitem{rosenberg2022} E.~Rosenberg, S.~Gratton, G.~Efstathiou, MNRAS \textbf{517}, 4620
  (2022), arXiv:2205.10869.
\bibitem{carron2022} J.~Carron, M.~Mirmelstein, A.~Lewis, JCAP \textbf{09}, 039 (2022),
  arXiv:2206.07773.
\bibitem{louis2025} T.~Louis et al.\ (ACT), \emph{The Atacama Cosmology Telescope: DR6 Power
  Spectra, Likelihoods and $\Lambda$CDM Parameters}, arXiv:2503.14452 (2025).
\bibitem{madhavacheril2024} M.~S.~Madhavacheril et al.\ (ACT), ApJ \textbf{962}, 113 (2024),
  arXiv:2304.05203.
\bibitem{qu2024} F.~J.~Qu et al.\ (ACT), ApJ \textbf{962}, 112 (2024), arXiv:2304.05202.
\bibitem{akaike1974} H.~Akaike, IEEE Trans.\ Autom.\ Control \textbf{19}, 716 (1974).
\bibitem{luminet2003} J.-P.~Luminet, J.~R.~Weeks, A.~Riazuelo, R.~Lehoucq, J.-P.~Uzan,
  Nature \textbf{425}, 593 (2003), arXiv:astro-ph/0310253.
\bibitem{planck2016topo} Planck Collaboration, A\&A \textbf{594}, A18 (2016), arXiv:1502.01593.
\bibitem{domagala2004} M.~Domagala, J.~Lewandowski, Class.\ Quantum Grav.\ \textbf{21}, 5233
  (2004), arXiv:gr-qc/0407051.
\bibitem{meissner2004} K.~A.~Meissner, Class.\ Quantum Grav.\ \textbf{21}, 5245 (2004),
  arXiv:gr-qc/0407052.
\bibitem{quevillon2019} J.~Quevillon, C.~Smith, Eur.\ Phys.\ J.\ C \textbf{79}, 822 (2019),
  arXiv:1903.12559.
\bibitem{carroll1998} S.~M.~Carroll, Phys.\ Rev.\ Lett.\ \textbf{81}, 3067 (1998),
  arXiv:astro-ph/9806099.
\bibitem{murai2023} K.~Murai, F.~Naokawa, T.~Namikawa, E.~Komatsu, Phys.\ Rev.\ D
  \textbf{107}, L041302 (2023), arXiv:2209.07804.
\bibitem{minami2020} Y.~Minami, E.~Komatsu, Phys.\ Rev.\ Lett.\ \textbf{125}, 221301 (2020),
  arXiv:2011.11254.
\bibitem{eskilt2022} J.~R.~Eskilt, E.~Komatsu, Phys.\ Rev.\ D \textbf{106}, 063503 (2022),
  arXiv:2205.13962.
\bibitem{so2019} P.~Ade et al.\ (Simons Observatory), JCAP \textbf{02}, 056 (2019),
  arXiv:1808.07445.
\bibitem{litebird2023} LiteBIRD Collaboration, Prog.\ Theor.\ Exp.\ Phys.\ \textbf{2023},
  042F01 (2023), arXiv:2202.02773.
\bibitem{kostelecky2008} V.~A.~Kosteleck\'y, N.~Russell, J.~D.~Tasson, Phys.\ Rev.\ Lett.\
  \textbf{100}, 111102 (2008), arXiv:0712.4393.
\end{thebibliography}

\end{document}
